\documentclass[11pt,letterpaper]{article}

\usepackage[margin=1in]{geometry}
\usepackage{parskip}
\usepackage{xcolor}
\usepackage{listings}
\usepackage{hyperref}
\usepackage{enumitem}

\hypersetup{
  colorlinks=true,
  linkcolor=blue!60!black,
  urlcolor=blue!60!black,
  citecolor=blue!60!black
}

% ---------- listings style for PowerLoom / STELLA ----------
\definecolor{plkeyword}{RGB}{0,80,160}
\definecolor{plcomment}{RGB}{100,120,100}
\definecolor{plstring}{RGB}{160,40,40}
\definecolor{plprompt}{RGB}{80,80,80}

\lstdefinelanguage{PowerLoom}{
  morekeywords={defmodule,defconcept,defrelation,deffunction,defrule,
    assert,retract,ask,retrieve,in-module,in-package,in-dialect,
    clear-module,reset-features,documentation,forall,exists,and,or,not,
    implies,subset-of,subrelation-of,has-parent,has-child,true,false,
    powerloom,load,demo,help,stop,save-kb,load-kb,print-object,
    forward-chain,run-forward-rules},
  sensitive=false,
  morecomment=[l]{;},
  morestring=[b]",
  keywordstyle=\color{plkeyword}\bfseries,
  commentstyle=\color{plcomment}\itshape,
  stringstyle=\color{plstring}
}

\lstset{
  language=PowerLoom,
  basicstyle=\ttfamily\small,
  showstringspaces=false,
  columns=fullflexible,
  breaklines=true,
  frame=single,
  framerule=0.3pt,
  rulecolor=\color{black!30},
  xleftmargin=0.5em,
  xrightmargin=0.5em,
  aboveskip=0.8em,
  belowskip=0.8em
}

\lstdefinestyle{shell}{
  language=bash,
  keywordstyle=\color{plkeyword}\bfseries,
  morekeywords={cd,sbcl}
}

\newcommand{\cmd}[1]{\texttt{#1}}
\newcommand{\file}[1]{\texttt{#1}}
\newcommand{\prompt}[1]{\textcolor{plprompt}{\texttt{#1}}}

% ---------- metadata ----------
\title{A Step-by-Step PowerLoom Tutorial}
\author{For use with PowerLoom running under SBCL}
\date{}

\begin{document}
\maketitle

\begin{abstract}
\noindent
This tutorial walks a newcomer from a cold terminal to a working PowerLoom
knowledge base. Every step is concrete: type the command, see the expected
response, understand what just happened. By the end you will have defined
modules, concepts, and relations; asserted and retrieved facts; written a
rule; performed classification; and saved your work. The tutorial assumes
no prior PowerLoom or STELLA experience, only a little comfort at a Lisp
prompt.
\end{abstract}

\tableofcontents
\bigskip\hrule

% ==========================================================
\section{Prerequisites}
\label{sec:prereqs}

Before starting, you need a PowerLoom that launches under SBCL. If you have
not already got that working, follow \file{../StartUp.md} in the PowerLoom
project directory --- it is a short, verified cold-start guide. The one-line
summary is:

\begin{lstlisting}[style=shell]
cd <path-to>/PowerLoom
sbcl --load load-powerloom.lisp
\end{lstlisting}

You should land at the SBCL prompt with the message
\cmd{PowerLoom <version> loaded.}

Throughout this tutorial:
\begin{itemize}[nosep]
  \item \prompt{*} marks the SBCL top-level prompt.
  \item \prompt{PL-USER |=} marks the PowerLoom listener prompt.
  \item Lines beginning with \cmd{;} are PowerLoom comments --- safe to copy.
\end{itemize}

% ==========================================================
\section{Starting and Stopping the Listener}
\label{sec:listener}

From the SBCL \prompt{*} prompt, enter the PowerLoom listener:

\begin{lstlisting}
(powerloom)
\end{lstlisting}

\noindent The prompt changes to \prompt{PL-USER |=}. You are now inside the
PowerLoom command loop, inside the default user module \cmd{PL-USER}. Try:

\begin{lstlisting}
(help)
\end{lstlisting}

\noindent PowerLoom prints a list of its top-level commands. To leave the
listener and return to SBCL:

\begin{lstlisting}
(stop)
\end{lstlisting}

\noindent You are back at the \prompt{*} prompt, but PowerLoom remains
loaded in memory --- so \cmd{(powerloom)} re-enters instantly. Use
\cmd{(sb-ext:exit)} (or \cmd{Ctrl-D}) to quit SBCL entirely when you are
done.

For the rest of this tutorial, re-enter the listener and stay there unless
told otherwise.

% ==========================================================
\section{Your First Module}
\label{sec:modules}

A \textbf{module} is a PowerLoom namespace: it holds concept, relation, and
instance definitions, and it holds the assertions you make against them.
Modules form a hierarchy. The out-of-the-box module is \cmd{PL-USER}, which
inherits from \cmd{PL-KERNEL} (where the built-ins live).

Before doing any real work, create a fresh module of your own so that this
tutorial does not collide with PL-USER. At the \prompt{PL-USER |=} prompt:

\begin{lstlisting}
(defmodule "TUTORIAL"
  :includes ("PL-USER"))
\end{lstlisting}

\noindent PowerLoom responds with something like \cmd{|MDL|TUTORIAL}. The
\cmd{|MDL|} prefix just tells you ``this is a module object.''

Now step into it:

\begin{lstlisting}
(in-module "TUTORIAL")
\end{lstlisting}

\noindent The prompt becomes \prompt{TUTORIAL |=}. Everything you define
from here lands in \cmd{TUTORIAL}, and \cmd{TUTORIAL} can see everything in
\cmd{PL-USER} and \cmd{PL-KERNEL}.

\textit{Tip:} if you want to start over inside a module without restarting
PowerLoom, use \cmd{(clear-module "TUTORIAL")} followed by the definitions
you want to keep.

% ==========================================================
\section{Concepts and Relations}
\label{sec:concepts}

A \textbf{concept} is a named class of individuals. A \textbf{relation}
connects individuals. We will model a small piece of a company ontology.

\subsection{Defining concepts}

\begin{lstlisting}
(defconcept Person
  :documentation "The class of persons.")

(defconcept Employee (Person)
  :documentation "A person who is employed somewhere.")

(defconcept Manager (Employee)
  :documentation "An employee who manages other employees.")

(defconcept Company
  :documentation "An organization that employs people.")
\end{lstlisting}

\noindent Each \cmd{defconcept} returns the concept object, e.g.\ \cmd{|C|Employee}.
The parenthesised argument --- \cmd{(Person)} for \cmd{Employee} --- names
parent concepts. \cmd{Employee} is therefore a subclass of \cmd{Person}.

\subsection{Defining relations}

\begin{lstlisting}
(defrelation works-for ((?e Employee) (?c Company))
  :documentation "True when ?e is employed by ?c.")

(defrelation manages ((?m Manager) (?e Employee))
  :documentation "True when ?m manages ?e.")

(defrelation salary ((?e Employee) (?s INTEGER))
  :documentation "?e's annual salary in dollars.")
\end{lstlisting}

\noindent Variables are written with a \cmd{?} prefix. The type after each
variable declares the expected argument sort. \cmd{INTEGER} is one of the
built-in STELLA types available in every module.

% ==========================================================
\section{Asserting Facts}
\label{sec:assert}

A fresh module has structure but no facts. Add some:

\begin{lstlisting}
(assert (Company Acme))
(assert (Person  Alice))
(assert (Person  Bob))
(assert (Person  Carol))

(assert (Employee Alice))
(assert (Employee Bob))
(assert (Manager  Carol))   ; Carol is automatically an Employee and a Person

(assert (works-for Alice Acme))
(assert (works-for Bob   Acme))
(assert (works-for Carol Acme))

(assert (manages Carol Alice))
(assert (manages Carol Bob))

(assert (= (salary Alice) 75000))
(assert (= (salary Bob)   80000))
(assert (= (salary Carol) 120000))
\end{lstlisting}

\noindent Every \cmd{assert} returns the proposition it just installed ---
PowerLoom's way of confirming. If you assert something structurally
inconsistent (wrong arity, wrong type), PowerLoom complains immediately.

% ==========================================================
\section{Asking Questions}
\label{sec:ask}

PowerLoom distinguishes two query styles:

\begin{description}[nosep,style=nextline]
  \item[\cmd{ask}] Yes/no. Returns \cmd{TRUE}, \cmd{FALSE}, or
    \cmd{UNKNOWN}.
  \item[\cmd{retrieve}] Produces bindings for free variables ---
    PowerLoom's ``SELECT.''
\end{description}

\subsection{Yes/no queries}

\begin{lstlisting}
(ask (Employee Alice))            ; TRUE (Alice was asserted Employee)
(ask (Employee Carol))            ; TRUE (via Manager -> Employee)
(ask (Person Carol))              ; TRUE (via Employee -> Person)
(ask (works-for Alice Acme))      ; TRUE
(ask (manages Alice Carol))       ; FALSE (inverted)
(ask (Manager Bob))               ; UNKNOWN -- never asserted or derivable
\end{lstlisting}

\noindent Notice how Carol's manager-ness cascades: \cmd{(Manager Carol)}
was the only explicit class assertion, yet \cmd{(Person Carol)} holds
because concept hierarchy is understood by the reasoner.

\subsection{Bindings queries}

\cmd{retrieve} takes a pattern with free variables and returns all matches.

\begin{lstlisting}
(retrieve all ?x (Employee ?x))
\end{lstlisting}

\noindent Prints each employee. The word \cmd{all} asks for every match; you
can also write \cmd{(retrieve 5 ?x ...)} to cap at five, or leave it out
for a single answer.

Multi-variable retrieval looks natural:

\begin{lstlisting}
(retrieve all (?m ?e)
  (manages ?m ?e))
\end{lstlisting}

\noindent Conjunctions and arithmetic work in retrieval too:

\begin{lstlisting}
(retrieve all ?e
  (and (Employee ?e)
       (> (salary ?e) 78000)))
\end{lstlisting}

% ==========================================================
\section{Retraction}
\label{sec:retract}

Retraction removes a previously asserted fact:

\begin{lstlisting}
(retract (manages Carol Bob))
(ask (manages Carol Bob))          ; FALSE (or UNKNOWN)
(assert (manages Carol Bob))       ; put it back
\end{lstlisting}

\noindent Retraction removes only the direct assertion. Facts that can still
be derived by inference remain true.

% ==========================================================
\section{Your First Rule}
\label{sec:rules}

A rule is an implication the reasoner applies on demand. Say we want every
employee whose salary exceeds \$100{,}000 to be classified as a
\cmd{HighEarner}:

\begin{lstlisting}
(defconcept HighEarner (Employee))

(defrule high-earner-rule
  (=> (and (Employee ?e) (> (salary ?e) 100000))
      (HighEarner ?e)))
\end{lstlisting}

\noindent Verify the rule fires:

\begin{lstlisting}
(ask (HighEarner Carol))           ; TRUE
(retrieve all ?x (HighEarner ?x))
\end{lstlisting}

\noindent PowerLoom's default reasoner is a backward chainer: it runs your
rule only when a query needs it. If you prefer to push consequences
eagerly, use \cmd{(run-forward-rules)}.

% ==========================================================
\section{Classification: A Family Ontology}
\label{sec:classify}

A signature PowerLoom feature is automatic classification. A concept
defined with an ``if and only if'' body lets the reasoner decide
membership for you. Clear the module and try the canonical family
example:

\begin{lstlisting}
(clear-module "TUTORIAL")
(in-module "TUTORIAL")

(defconcept Human  :documentation "Human beings.")
(defconcept Woman (Human))
(defconcept Man   (Human))

(defrelation has-child ((?p Human) (?c Human))
  :inverse has-parent)

(defconcept Mother (?m Woman)
  :<=> (exists (?x Human) (has-child ?m ?x)))

(defconcept Father (?f Man)
  :<=> (exists (?x Human) (has-child ?f ?x)))

(defconcept Parent (?p Human)
  :<=> (or (Mother ?p) (Father ?p)))

(defconcept Grandmother (?g Woman)
  :<=> (exists (?c Parent) (has-child ?g ?c)))
\end{lstlisting}

\noindent Assert facts:

\begin{lstlisting}
(assert (Woman Elizabeth))
(assert (Man   Charles))
(assert (Man   William))
(assert (has-child Elizabeth Charles))
(assert (has-child Charles   William))
\end{lstlisting}

\noindent Now query \emph{derived} concept membership --- nothing below was
ever asserted explicitly:

\begin{lstlisting}
(ask (Mother      Elizabeth))      ; TRUE
(ask (Grandmother Elizabeth))      ; TRUE
(ask (Father      Charles))        ; TRUE
(ask (Parent      Charles))        ; TRUE
\end{lstlisting}

\noindent The reasoner also proves subclass relationships between concepts
themselves, using just the definitions:

\begin{lstlisting}
(ask (subset-of Grandmother Mother))   ; TRUE
(ask (subset-of Grandmother Parent))   ; TRUE
\end{lstlisting}

\noindent This is \emph{classification} --- arguably PowerLoom's most
distinctive trick.

% ==========================================================
\section{Saving and Loading Your KB}
\label{sec:persist}

Edit and persist your work in a file rather than retyping into the
listener. Create a file, say \file{my-company.plm}, with content like:

\begin{lstlisting}
(in-package "STELLA")

(defmodule "MY-COMPANY"
  :includes ("PL-USER"))

(in-module "MY-COMPANY")

(defconcept Employee)
(defrelation works-for ((?e Employee) (?c OBJECT)))

(assert (Employee Alice))
(assert (works-for Alice Acme))
\end{lstlisting}

\noindent Then load it from the PowerLoom listener:

\begin{lstlisting}
(load "my-company.plm")
(in-module "MY-COMPANY")
(retrieve all ?x (Employee ?x))
\end{lstlisting}

\noindent To dump a module's current state to a file programmatically:

\begin{lstlisting}
(save-kb "snapshot.plm" "MY-COMPANY")
\end{lstlisting}

\noindent \cmd{save-kb} writes every definition and assertion in the named
module; you can \cmd{(load "snapshot.plm")} in a future session to get
back to exactly this state.

% ==========================================================
\section{Exploring the Bundled Demos}
\label{sec:demos}

Twenty-plus prepackaged demos ship with PowerLoom in
\file{sources/logic/demos/}. From the listener, launch the picker:

\begin{lstlisting}
(demo)
\end{lstlisting}

\noindent Recommended order for a new user:

\begin{enumerate}[nosep]
  \item \file{basics.plm} --- modules, definitions, retrieval, retraction.
  \item \file{family.plm} --- the classification example from
        Section~\ref{sec:classify}, in more depth.
  \item \file{relation-hierarchy.plm} --- subrelations and inheritance.
  \item \file{constraints.plm} --- using rules to enforce invariants.
  \item \file{kinship.plm} --- a larger family knowledge base.
  \item \file{aircraft.plm} --- a realistic classification ontology.
\end{enumerate}

\noindent Each demo file is heavily commented; reading them is as valuable
as running them.

% ==========================================================
\section{Where to Go Next}
\label{sec:next}

\begin{itemize}[nosep]
  \item \textbf{The manual.} \file{sources/logic/doc/} in the PowerLoom
        tree contains the reference manual in PDF and HTML.
  \item \textbf{Course materials.} \file{PowerLoom-Course-Lecture.pdf}
        and \file{PowerLoom-Overview-Slides.pdf} in this same
        \file{tutorials/} directory give a lecture-paced introduction.
  \item \textbf{The PLI API.} Embedding PowerLoom in a host application
        (Lisp, C++, Java, or Python) is done through the PowerLoom
        Language Interface; see \file{sources/logic/pli.ste} and the
        Python bindings under \file{python/}.
  \item \textbf{Extension systems.} RDBMS access, the web UI
        (Ontosaurus), and others can be enabled by setting
        \cmd{*load-all-extensions?*} to \cmd{t} before loading
        PowerLoom (see \file{../StartUp.md} step~8).
\end{itemize}

% ==========================================================
\section*{Appendix: Building This Document}

This tutorial compiles with any standard LaTeX distribution:

\begin{lstlisting}[style=shell]
cd <path-to>/PowerLoom/tutorials
pdflatex PowerLoom-Tutorial.tex
pdflatex PowerLoom-Tutorial.tex     # second pass for the TOC
\end{lstlisting}

\noindent The output is \file{PowerLoom-Tutorial.pdf}.

\end{document}
